\chapter{扫描与并行算法的工作效率}
\label{chap:scan}

\setcounter{footnote}{0}

本章特别感谢 Li-Wen Chang、Juan G\'omez-Luna 和 John Owens 的贡献。

并行扫描（parallel scan）经常用来并行化表面上顺序的操作，例如资源分配、工作分派和
多项式求值。一般而言，如果一项计算天然写成数学递推式，即序列中的每一项都由前一项
定义，那么它往往可以改写成并行扫描。并行扫描在大规模并行计算中十分关键，原因很
简单：应用中任何顺序部分都可能严重限制整体性能，而其中许多部分都能借助扫描变成并行
计算。因此，过滤、基数排序、快速排序、字符串比较、多项式求值、递推式求解和树操作等
并行算法，常把扫描作为基本操作。第 12 章将介绍过滤，第 14 章将介绍基数排序。

并行扫描算法不止一种，它们在工作复杂度与跨度之间有不同取舍。有些并行算法完成的工作
多于顺序算法；硬件资源受限时，增加的工作甚至会让并行扫描比顺序扫描更慢，因而必须用
优化来缓解。本章以扫描为例介绍双缓冲和单向同步等新技术，并再次运用 warp 级原语、线程
粗化和寄存器分块等已经见过的优化。

\section{背景}
\label{sec:scan-background}

数学上，包含式扫描（inclusive scan）接收二元运算符 $\oplus$ 和包含 $N$ 个元素的输入
数组 $[x_0,x_1,\ldots,x_{N-1}]$，返回
\[
 [x_0,(x_0\oplus x_1),\ldots,(x_0\oplus x_1\oplus\cdots\oplus x_{N-1})].
\]
例如，当 $\oplus$ 为加法时，对 $[3,1,7,0,4,1,6,3]$ 扫描得到
$[3,4,11,11,15,16,22,25]$。之所以称为“包含式”，是因为每个输出元素包含对应输入
元素的贡献。以加法为运算符的扫描也常称为\textbf{前缀和（prefix sum）}。

可以用为一群人切香肠的例子说明扫描。假设要把一根 40 英寸的香肠分给 8 个人，各人依次
要 $3,1,7,0,4,1,6,3$ 英寸。顺序做法是先为第 0 人切 3 英寸，剩下 37 英寸；再为第 1
人切 1 英寸，剩下 36 英寸；如此继续，最后为第 7 人切下 3 英寸。总共切出 25 英寸，
剩余 15 英寸。

包含式扫描则直接根据每个人所需长度算出所有切点。输入
$[3,1,7,0,4,1,6,3]$ 的前缀和 $[3,4,11,11,15,16,22,25]$ 恰好就是切点位置。
第一个切点在 3 英寸处，得到第 0 人所需的 3 英寸；第二个在 4 英寸处，两点之间正好
1 英寸；最后一个在 25 英寸处，与前一切点 22 英寸相距 3 英寸。既然所有切点都已知，
八刀就可以同时落下，也可以按任意次序执行。

直观地说，包含式扫描接收一组人的请求，找出能一次满足所有请求的分界点。请求对象可以
是香肠、面包、露营地，也可以是计算机中的一段连续内存；只要快速算出全部分界点，就能
并行满足各项请求。

排除式扫描（exclusive scan）与包含式扫描相似，但输出排列略有不同：
\[
 [\mathrm{ID}_{\oplus},x_0,(x_0\oplus x_1),\ldots,
  (x_0\oplus x_1\oplus\cdots\oplus x_{N-2})].
\]
每个输出元素都排除了对应输入元素的贡献。首元素是运算符 $\oplus$ 的单位元
$\mathrm{ID}_{\oplus}$，末元素只包含到 $x_{N-2}$ 为止的贡献。二元运算符的单位元作为
一个操作数时，结果等于另一操作数；加法的单位元是 0。

排除式扫描的用途与包含式扫描基本相同，只是提供的信息略有不同。香肠例中，它返回
$[0,3,4,11,11,15,16,22]$，这些值是各段的起点。第 0 人的一段从 0 英寸处开始，
第 7 人的一段从 22 英寸处开始。内存分配等应用需要把已分配区域的起始指针返回给请求
者，因此起点信息尤其重要。

两种扫描结果很容易互相转换。包含式结果整体右移一位并在第 0 项填入单位元，就得到
排除式结果；排除式结果整体左移一位，再把“原末元素 $\oplus$ 最后一个输入元素”填入
末项，就得到包含式结果。直接生成哪一种只是使用切点还是区段起点的便利选择。因此，
本章后续只给出包含式扫描的并行算法和实现。

在讨论并行算法之前，先看图 \ref{fig:11-1} 的顺序包含式扫描。假设运算符为加法，
输入和输出分别放在 \code{input} 与 \code{output} 中。代码先以
\code{input[0]} 初始化 \code{output[0]}，随后每次循环把一个新输入加到保存此前
所有输入累积值的前一输出上。对 $N$ 个元素，它执行 $N-1$ 次加法，因此计算复杂度为
$O(N)$。

\begin{figure}[htbp]
  \centering
  \begin{minipage}{0.9\textwidth}
  \begin{lstlisting}[language=C++,numbers=left]
void sequential_scan(float *input, float *output, unsigned int N) {
    output[0] = input[0];
    for (unsigned int i = 1; i < N; ++i) {
        output[i] = output[i - 1] + input[i];
    }
}
  \end{lstlisting}
  \end{minipage}
  \caption{基于加法的简单顺序包含式扫描实现}
  \label{fig:11-1}
\end{figure}

\section{使用 Kogge--Stone 算法进行并行扫描}
\label{sec:kogge-stone-scan}

一种容易想到的并行化方法是：每个输出元素交给一个线程，该线程顺序归约相关的输入。
虽然所有输出可同时计算，但最后一个线程计算 $y_{N-1}$ 仍需 $N$ 个步，程序的完成时间
受最慢线程限制，所以跨度仍为 $O(N)$，并不优于顺序扫描。资源有限时，它甚至可能慢得多。
因为第 $i$ 个输出需要 $i$ 次计算，所有线程完成的运算数为
\[
  \sum_{i=0}^{N-1}i=\frac{N(N-1)}{2},
\]
即工作复杂度为 $O(N^2)$：增加了大量工作，却没有降低跨度。

更好的方法是用第 10 章的并行归约树计算各个输出。扫描运算符必须满足结合律；视优化
方式而定，有时还需要满足交换律。如果各输出独立建立归约树，工作量仍为 $O(N^2)$；
关键是让不同归约树共享部分和。Kogge--Stone 算法正是这样做的。它最初在 20 世纪 70
年代为高速加法器电路提出\cite{kogge1973parallel}，至今仍用于高速算术硬件。

图 \ref{fig:11-2} 展示一个含 16 个元素的原地扫描。开始时位置 $i$ 保存 $x_i$；完成
$k$ 轮后，它保存以该位置为末端、最多含 $2^k$ 个输入的和。第一轮把相邻元素相加，
第二轮把相距 2 的部分和相加，随后距离依次为 4 和 8。位置 15 的粗线依赖路径形成一棵
归约树。

\begin{figure}[htbp]
  \centering\scriptsize
  \begin{tabular}{c|l}
    初始 & $x_0\quad x_1\quad x_2\quad x_3\quad\cdots\quad x_{15}$\\
    距离 1 & $x_0\quad \sum_{0}^{1}x\quad \sum_{1}^{2}x\quad
      \sum_{2}^{3}x\quad\cdots\quad\sum_{14}^{15}x$\\
    距离 2 & $x_0\quad \sum_{0}^{1}x\quad \sum_{0}^{2}x\quad
      \sum_{0}^{3}x\quad\cdots\quad\sum_{12}^{15}x$\\
    距离 4 & 前 8 个位置逐步完成，后部每项合并左侧相距 4 的部分和\\
    距离 8 & 所有位置成为 $y_i=\sum_{j=0}^{i}x_j$
  \end{tabular}
  \caption{基于 Kogge--Stone 加法器设计的并行包含式扫描}
  \label{fig:11-2}
\end{figure}

第 0 个位置从开始就是最终答案 $y_0=x_0$。第一轮中，除位置 0 外，每个位置把当前值
与左邻相加，于是位置 $i$ 保存 $x_{i-1}+x_i$；位置 1 此时已经完成。第二轮中，除位置
0、1 外，每个位置再加上左侧相距 2 的值，位置 $i$ 于是保存
$x_{i-3}+x_{i-2}+x_{i-1}+x_i$；位置 2 和 3 此时完成。后续各轮同理。

图 \ref{fig:11-3} 实现每块局部扫描。块区段大小是编译期常量
\code{SEG_SIZE}，启动时块大小也取 \code{SEG_SIZE}，每个线程演化一个元素。
线程先协作把全局内存区段装入共享内存；越界位置填 0。循环的步幅每轮加倍，活跃线程
把当前位置与左侧相距 \code{stride} 的位置相加。末尾再把结果写回全局内存。

\begin{figure}[htbp]
  \centering
  \begin{minipage}{0.94\textwidth}
  \begin{lstlisting}[language=C++,numbers=left]
__global__ void scan_kernel(float *input, float *output, unsigned int N) {
    __shared__ float buffer_s[SEG_SIZE];
    unsigned int i = blockIdx.x * blockDim.x + threadIdx.x;
    if (i < N) buffer_s[threadIdx.x] = input[i];
    else       buffer_s[threadIdx.x] = 0.0f;

    for (unsigned int stride = 1; stride < blockDim.x; stride *= 2) {
        __syncthreads();
        float temp;
        if (threadIdx.x >= stride)
            temp = buffer_s[threadIdx.x] + buffer_s[threadIdx.x - stride];
        __syncthreads();
        if (threadIdx.x >= stride)
            buffer_s[threadIdx.x] = temp;
    }
    if (i < N) output[i] = buffer_s[threadIdx.x];
}
  \end{lstlisting}
  \end{minipage}
  \caption{Kogge--Stone 块内包含式扫描内核}
  \label{fig:11-3}
\end{figure}

循环每轮需要两次屏障。第一次保证其他线程读取前，所有线程已写完上一轮结果；第二次
保证任何线程覆盖共享内存前，其他线程已读完旧值。条件造成的控制分歧较轻：步幅小于
warp 大小时，只有第一个 warp 中的一部分线程退出；步幅达到 warp 大小时，退出的是整个
warp，不再产生 warp 内分歧。

第二个屏障是扫描区别于第 10 章归约内核的关键。每个活跃线程读取自己的位置和左侧
位置。如果线程 $i$ 在线程 $i+\mathrm{stride}$ 读取旧值前便覆盖自己的位置，后者会
错误地读到本轮新值。例如第二轮中线程 6 需要旧的
\code{buffer_s[4]}=$x_3+x_4$；若线程 4 先写入
$x_1+x_2+x_3+x_4$，线程 6 会多计 $x_1+x_2$，下一轮又会重复计入，最终结果错误。
结果是否正确将依赖线程时序，甚至不同运行之间也不一致。

临时变量 \code{temp} 与第二次 \code{__syncthreads()} 共同消除竞态：所有活跃
线程先把加法结果写入私有变量，屏障确认旧值全部读完后，才覆盖共享内存。归约树无需
第二个屏障，是因为每轮被活跃线程写入的位置不会在同一轮被其他活跃线程读取。扫描中
并发归约树共享部分和，写入位置可能恰好是别人的输入，因此必须同步。

\section{用双缓冲减少同步}
\label{sec:scan-double-buffering}

图 \ref{fig:11-3} 的主要瓶颈是屏障开销。每轮两次屏障所约束的依赖不同：第一次保证
上一轮写完后本轮才读，是真依赖；第二次保证本轮读完后才覆盖，是假依赖。第 6 章已经
说明，可以用双缓冲消除约束假依赖的屏障。

一般地，把每轮输入与输出存到不同缓冲区，就不会覆盖同轮还要读取的数据。为每轮各分配
一个缓冲区会消耗大量共享内存；但一轮输出算完后，该轮输入就不再需要，所以只要两个
缓冲区轮流担任输入和输出即可。

\begin{figure}[htbp]
  \centering\small
  \begin{tabular}{c|c|c}
    阶段 & 输入缓冲区 & 输出缓冲区\\\hline
    装载 & 全局内存 $\rightarrow$ buffer1 & ---\\
    轮 0（距离 1） & buffer1 & buffer2\\
    轮 1（距离 2） & buffer2 & buffer1\\
    轮 2（距离 4） & buffer1 & buffer2\\
    后续 & \multicolumn{2}{c}{每轮交换两个指针}
  \end{tabular}
  \caption{把双缓冲应用于 Kogge--Stone 并行扫描}
  \label{fig:11-4}
\end{figure}

双缓冲后，已完成、不再演化的值不会自动留在输出缓冲区，必须显式复制。例如轮 0 中
线程 0 不做加法，但要把 $x_0$ 从 \code{buffer1} 复制到 \code{buffer2}。

\begin{figure}[htbp]
  \centering
  \begin{minipage}{0.95\textwidth}
  \begin{lstlisting}[language=C++,numbers=left]
__global__ void scan_kernel(float *input, float *output, unsigned int N) {
    unsigned int i = blockIdx.x * blockDim.x + threadIdx.x;
    __shared__ float buffer1_s[SEG_SIZE];
    __shared__ float buffer2_s[SEG_SIZE];
    float *inBuffer_s  = buffer1_s;
    float *outBuffer_s = buffer2_s;
    inBuffer_s[threadIdx.x] = (i < N) ? input[i] : 0.0f;

    for (unsigned int stride = 1; stride < blockDim.x; stride *= 2) {
        __syncthreads();
        if (threadIdx.x >= stride) {
            outBuffer_s[threadIdx.x] = inBuffer_s[threadIdx.x]
                                      + inBuffer_s[threadIdx.x - stride];
        } else {
            outBuffer_s[threadIdx.x] = inBuffer_s[threadIdx.x];
        }
        float *temp = inBuffer_s;
        inBuffer_s = outBuffer_s;
        outBuffer_s = temp;
    }
    if (i < N) output[i] = inBuffer_s[threadIdx.x];
}
  \end{lstlisting}
  \end{minipage}
  \caption{采用双缓冲的 Kogge--Stone 块内包含式扫描内核}
  \label{fig:11-5}
\end{figure}

图 \ref{fig:11-5} 用两个共享内存数组和两个指针交换角色。参与加法的线程从输入缓冲
读取、向输出缓冲写入，其他线程执行复制，每轮末尾交换指针。循环只剩一次屏障，用来
维持真依赖；所有读来自 \code{inBuffer_s}，所有写进入
\code{outBuffer_s}，读后写的假依赖已经消失。

\section{用 warp 级原语减少同步}
\label{sec:scan-warp-primitives}

双缓冲虽然减半了屏障数，开销仍然可观：图 \ref{fig:11-5} 每轮只完成一次浮点加法，
却调用一次 \code{__syncthreads()} 并访问三次共享内存，循环明显受同步与共享内存
访问主导。warp 内线程可以用 warp 级原语更高效地同步和交换数据。因此，可以让每个
warp 先扫描自己的子区段，再合并各 warp 的结果，得到整个块的扫描结果。

一种分解方式是\textbf{扫描--扫描--加法（scan--scan--add）}：先分别扫描较小区段，
再扫描各区段的和，最后把此前区段的累计和加到当前区段。图 \ref{fig:11-6} 以 16 个
元素分成 4 段为例。第一阶段得到每段内部的扫描结果；各段末元素 7、7、6、11 是段和。
第二阶段扫描段和得到 7、14、20、31，它们分别是原问题在位置 3、7、11、15 的最终
结果。第三阶段把 7 加到第二段、14 加到第三段、20 加到第四段，补齐所有结果。

\begin{figure}[htbp]
  \centering\small
  \begin{tabular}{c|c|c|c|c}
    & 段 0 & 段 1 & 段 2 & 段 3\\\hline
    输入 & 2,1,3,1 & 0,4,1,2 & 0,3,1,2 & 4,1,2,4\\
    段内扫描 & 2,3,6,7 & 0,4,5,7 & 0,3,4,6 & 4,5,7,11\\
    段和扫描 & \multicolumn{4}{c}{7,14,20,31}\\
    加前缀后 & 2,3,6,7 & 7,11,12,14 & 14,17,18,20 & 24,25,27,31
  \end{tabular}
  \caption{扫描--扫描--加法分解示例}
  \label{fig:11-6}
\end{figure}

块级扫描可以同样分解到多个 warp：每个 warp 扫描自己的连续子区段；每个 warp 的最后
一个线程把 warp 和写入共享内存；第一个 warp 扫描这些和；最后，各 warp 把此前所有
warp 的总和加到自己的扫描结果上。

\begin{figure}[htbp]
  \centering
  \fbox{\begin{minipage}{0.88\textwidth}\centering
    块输入 $\longrightarrow$ 各 warp 独立扫描\\[0.4em]
    warp 末线程收集段和 $\longrightarrow$ warp 0 扫描段和\\[0.4em]
    每个 warp 加上前序 warp 的累计和 $\longrightarrow$ 块级扫描结果
  \end{minipage}}
  \caption{用扫描--扫描--加法把块级扫描分解为 warp 级扫描}
  \label{fig:11-7}
\end{figure}

图 \ref{fig:11-8} 的设备函数用 Kogge--Stone 方法完成 warp 级扫描。每轮通过
\texttt{\_\_shfl\allowbreak\_up\allowbreak\_sync} 取得左侧相距 \code{stride} 的线程值，不需要共享内存或
块级屏障；只有 lane 下标不小于步幅的线程参与加法。

\begin{figure}[htbp]
  \centering
  \begin{minipage}{0.88\textwidth}
  \begin{lstlisting}[language=C++,numbers=left]
__device__ inline float warpScan(float val) {
    for (unsigned int stride = 1; stride < WARP_SIZE; stride *= 2) {
        float leftVal = __shfl_up_sync(0xffffffff, val, stride);
        if (laneIdx() >= stride) val += leftVal;
    }
    return val;
}
  \end{lstlisting}
  \end{minipage}
  \caption{使用 warp 级原语执行 warp 级扫描的设备函数}
  \label{fig:11-8}
\end{figure}

图 \ref{fig:11-9} 进一步组成块级扫描。各 warp 先独立扫描；最后一个 lane 写出 warp
和；第一个 warp 扫描这些和并写回共享内存；除第一个 warp 外，每个线程再加上前一个
warp 的扫描和。两个屏障分别位于第一次与第二次扫描之间、第二次扫描与最终加法之间。

\begin{figure}[htbp]
  \centering
  \begin{minipage}{0.92\textwidth}
  \begin{lstlisting}[language=C++,numbers=left]
__device__ inline float blockScan(float val) {
    val = warpScan(val);
    __shared__ float warpSums_s[NUM_WARPS];
    if (laneIdx() == WARP_SIZE - 1)
        warpSums_s[warpIdx()] = val;
    __syncthreads();

    if (warpIdx() == 0) {
        float warpSum = (threadIdx.x < NUM_WARPS)
                        ? warpSums_s[threadIdx.x] : 0.0f;
        warpSum = warpScan(warpSum);
        if (threadIdx.x < NUM_WARPS)
            warpSums_s[threadIdx.x] = warpSum;
    }
    __syncthreads();
    if (warpIdx() > 0) val += warpSums_s[warpIdx() - 1];
    return val;
}
  \end{lstlisting}
  \end{minipage}
  \caption{使用扫描--扫描--加法和 warp 级扫描执行块级扫描的设备函数}
  \label{fig:11-9}
\end{figure}

与图 \ref{fig:11-5} 每轮都访问共享内存并同步不同，图 \ref{fig:11-9} 只在两个阶段
边界处执行这些操作。于是，调用块级函数的内核非常简单：线程找出自己的元素，装入
寄存器，调用 \code{blockScan}，再把结果写回。

\begin{figure}[htbp]
  \centering
  \begin{minipage}{0.88\textwidth}
  \begin{lstlisting}[language=C++,numbers=left]
__global__ void scan_kernel(float *input, float *output, unsigned int N) {
    unsigned int i = blockIdx.x * blockDim.x + threadIdx.x;
    float val = (i < N) ? input[i] : 0.0f;
    val = blockScan(val);
    if (i < N) output[i] = val;
}
  \end{lstlisting}
  \end{minipage}
  \caption{调用 warp 级实现的块级扫描内核}
  \label{fig:11-10}
\end{figure}

\section{工作效率考量}
\label{sec:scan-work-efficiency}

分析并行算法时，一个重要指标是\textbf{工作效率（work efficiency）}，即算法实际完成
的工作量在多大程度上接近计算所需的最少工作量。\footnote{这里把“最少工作量”经验性
定义为所有已知算法中最小的工作量。}扫描 $N$ 个元素至少需要 $N-1$ 次加法，图
\ref{fig:11-1} 的顺序算法正好如此。最朴素的并行算法却完成 $N(N-1)/2$ 次加法，
即 $O(N^2)$，因而不是工作高效的。

分析图 \ref{fig:11-2} 的 Kogge--Stone 算法。扫描 $N$ 个元素需要 $\log_2N$ 轮。
步幅为 \code{stride} 的一轮有 $N-\mathrm{stride}$ 个活跃线程，每个线程完成一次
加法，所以总工作量是
\[
 \sum_{\mathrm{stride}\in\{1,2,4,\ldots,N/2\}}
 (N-\mathrm{stride})
 =N\log_2N-(N-1).
\]
真实 CUDA GPU 上，虽然越到后期退出的线程越多，只要所属 warp 尚未整体完成，它们仍
占用执行资源。因此，资源消耗更接近 $N\log_2N$，工作复杂度为 $O(N\log N)$。

好的一面是，这优于朴素并行方法的 $O(N^2)$；坏的一面是，仍不如顺序算法的 $O(N)$。
即使区段不大，额外工作也很显著：扫描 512 个元素时，并行算法完成的工作约为顺序代码
的 8 到 9 倍。

Kogge--Stone 的优势在于用并行执行减少步数。顺序循环执行 $N$ 轮，跨度为 $O(N)$；
并行内核最多执行 $\log_2N$ 轮，跨度为 $O(\log N)$。资源无限时，步数减少约
$N/\log_2N$ 倍；$N=512$ 时约为 $512/9=56.9$ 倍。

“步”可以近似比较扫描算法。顺序扫描 $N$ 个元素约需 $N$ 个步。CUDA 设备有 $P$ 个
执行单元时，Kogge--Stone 内核约需 $(N\log_2N)/P$ 个步。若 $P=N$，就是
$\log_2N$ 个步；若用 1024 个线程、32 个执行单元处理 1024 个元素，则约需
$(1024\times10)/32=320$ 个步，相对顺序执行的 1024 个步减少约 $3.2$ 倍。

额外工作有两方面问题。其一，硬件利用效率降低；$P$ 很小时，并行算法可能需要比顺序
算法更多的步，反而更慢。其二，即使运行更快，额外运算也会消耗额外能量，因此不适合
移动设备等功耗受限环境。

\section{用粗化提高工作效率}
\label{sec:scan-coarsening}

并行扫描的主要开销来自工作效率下降。执行资源充足时，以更多工作换取更短跨度是值得
的；资源受限、硬件不得不串行调度线程时，这些开销便没有必要。可以主动使用线程粗化，
把一部分工作串行化，从而降低并行化开销并提高工作效率。

具体做法是，让每个线程先对多个输入元素执行工作高效的顺序扫描，再由线程协作完成
并行扫描。扫描--扫描--加法分解同样适用：把块区段划分为每线程一个连续子区段；各线程
顺序扫描自己的子区段；用各子区段末值参与块级扫描；最后把此前所有线程的累计和加到
本线程的各个局部结果上。

\begin{figure}[htbp]
  \centering
  \fbox{\begin{minipage}{0.88\textwidth}\centering
    共享内存中的块区段\\[0.35em]
    $\Downarrow$ 每线程取得一个连续子区段并顺序扫描\\[0.35em]
    子区段末值 $\Downarrow$ 块级并行扫描线程和\\[0.35em]
    每线程加上前序线程累计和 $\Downarrow$ 完整块级扫描结果
  \end{minipage}}
  \caption{用扫描--扫描--加法把线程粗化应用于并行扫描}
  \label{fig:11-11}
\end{figure}

图 \ref{fig:11-12} 中，常量 \code{COARSE_FACTOR} 是每线程负责的元素数，
\code{BLOCK_DIM} 是块大小，所以每块区段包含
\code{COARSE_FACTOR*BLOCK_DIM} 个元素。线程虽然最终拥有连续子区段，却不能各自
直接成段地装载，否则 warp 内地址相隔 \code{COARSE_FACTOR}，无法合并。代码把
块区段视为 \code{COARSE_FACTOR} 个大小为 \code{BLOCK_DIM} 的连续块，逐块
合并装入共享内存；输出时采用同样方式合并写回。

\begin{figure}[htbp]
  \centering
  \begin{minipage}{0.98\textwidth}
  \begin{lstlisting}[language=C++,numbers=left]
__global__ void scan_kernel(float *input, float *output, unsigned int N) {
    unsigned int blockSegment = blockIdx.x * COARSE_FACTOR * BLOCK_DIM;
    __shared__ float buffer_s[COARSE_FACTOR * BLOCK_DIM];
    for (unsigned int c = 0; c < COARSE_FACTOR; ++c) {
        buffer_s[c * BLOCK_DIM + threadIdx.x] =
            input[blockSegment + c * BLOCK_DIM + threadIdx.x];
    }
    __syncthreads();

    unsigned int threadSegment = threadIdx.x * COARSE_FACTOR;
    for (unsigned int c = 1; c < COARSE_FACTOR; ++c)
        buffer_s[threadSegment + c] += buffer_s[threadSegment + c - 1];

    float threadSum = buffer_s[threadSegment + COARSE_FACTOR - 1];
    threadSum = blockScan(threadSum);
    __shared__ float threadSums[BLOCK_DIM];
    threadSums[threadIdx.x] = threadSum;
    __syncthreads();

    if (threadIdx.x > 0) {
        float prevPartialSum = threadSums[threadIdx.x - 1];
        for (unsigned int c = 0; c < COARSE_FACTOR; ++c)
            buffer_s[threadSegment + c] += prevPartialSum;
    }
    __syncthreads();
    for (unsigned int c = 0; c < COARSE_FACTOR; ++c) {
        output[blockSegment + c * BLOCK_DIM + threadIdx.x] =
            buffer_s[c * BLOCK_DIM + threadIdx.x];
    }
}
  \end{lstlisting}
  \end{minipage}
  \caption{用扫描--扫描--加法进行线程粗化的并行扫描内核}
  \label{fig:11-12}
\end{figure}

装载完毕后，各线程以 \code{threadIdx.x*COARSE_FACTOR} 找到自己的子区段并执行顺序
扫描。末值代表线程全部元素之和，传给图 \ref{fig:11-9} 的 \code{blockScan}。
扫描后的线程和写入共享内存，除第一个线程外，每个线程读取前一个线程的扫描和——也就是
此前所有线程元素的总和——并把它加到自己的全部局部结果上。最后同步并按合并方式写回。

设用 $P$ 个线程扫描 $N$ 个元素。每个线程顺序扫描 $N/P$ 个元素，共做
$P(N/P-1)=N-P$ 次操作；并行扫描 $P$ 个线程和约做 $P\log_2P$ 次操作；最后 $P-1$
个线程向各自 $N/P$ 个结果加前缀，共做 $(P-1)N/P=N-N/P$ 次操作。总计
\[
  2N+P\log_2P-P-\frac NP.
\]
当 $P$ 接近 $N$ 时，仍是 $O(N\log N)$；当 $P\ll N$ 时则接近 $N$，与顺序算法一样
为 $O(N)$。因此，在执行资源远少于元素数时，粗化能显著提高工作效率。

步数也随之改善。各线程的顺序扫描需 $N/P-1$ 个步，线程和扫描需 $\log_2P$ 个步，
最后加前缀需 $N/P$ 个步，总计
\[
  \frac{2N}{P}+\log_2P-1.
\]
不粗化时约为 $N\log_2N/P$。当 $P$ 接近 $N$ 时，二者均为 $O(\log N)$；当
$P\ll N$ 时，粗化后的步数接近 $2N/P$，明显更少。

粗化还有其他收益。块级并行扫描需要屏障并会产生少量控制分歧，而线程级顺序扫描不需要；
把更多操作移到线程内部，便减弱了同步、共享内存访问和分歧的相对影响。这与第 10 章粗化
归约内核的收益相同。

此外，粗化可以在保持块大小不变的同时扩大块区段，使网格中的线程块数按粗化因子减少。
第 \ref{sec:global-scan-consolidation} 节将看到，块数减少有利于合并各块区段。当然，也
可以保持区段大小不变而减少块内线程数，但前一种方式更适合全局扫描。

\section{用寄存器分块避免共享内存访问延迟}
\label{sec:scan-register-tiling}

图 \ref{fig:11-12} 把每线程子区段放在共享内存中，所以每次访问都要读写共享内存。
这些元素只由所属线程使用，让线程反复从共享内存取回是浪费。可以采用寄存器分块，把
子区段保存在寄存器中并反复就地访问。

\begin{figure}[htbp]
  \centering
  \begin{minipage}{0.98\textwidth}
  \begin{lstlisting}[language=C++,numbers=left]
__global__ void scan_kernel(float *input, float *output, unsigned int N) {
    unsigned int blockSegment = blockIdx.x * COARSE_FACTOR * BLOCK_DIM;
    __shared__ float buffer_s[COARSE_FACTOR * BLOCK_DIM];
    for (unsigned int c = 0; c < COARSE_FACTOR; ++c)
        buffer_s[c * BLOCK_DIM + threadIdx.x] =
            input[blockSegment + c * BLOCK_DIM + threadIdx.x];
    __syncthreads();

    unsigned int threadSegment = threadIdx.x * COARSE_FACTOR;
    float buffer_r[COARSE_FACTOR];
    buffer_r[0] = buffer_s[threadSegment];
#pragma unroll
    for (unsigned int c = 1; c < COARSE_FACTOR; ++c)
        buffer_r[c] = buffer_s[threadSegment + c] + buffer_r[c - 1];

    float threadSum = buffer_r[COARSE_FACTOR - 1];
    threadSum = blockScan(threadSum);
    __shared__ float threadSums[BLOCK_DIM];
    threadSums[threadIdx.x] = threadSum;
    __syncthreads();
    float prevPartialSum = (threadIdx.x > 0)
                           ? threadSums[threadIdx.x - 1] : 0.0f;
#pragma unroll
    for (unsigned int c = 0; c < COARSE_FACTOR; ++c)
        buffer_s[threadSegment + c] = buffer_r[c] + prevPartialSum;

    __syncthreads();
    for (unsigned int c = 0; c < COARSE_FACTOR; ++c)
        output[blockSegment + c * BLOCK_DIM + threadIdx.x] =
            buffer_s[c * BLOCK_DIM + threadIdx.x];
}
  \end{lstlisting}
  \end{minipage}
  \caption{对线程子区段采用寄存器分块的并行扫描内核}
  \label{fig:11-13}
\end{figure}

图 \ref{fig:11-13} 声明局部数组 \code{buffer_r} 保存线程子区段。第一次从共享内存
装入时就留在局部数组中，后续扫描与加前缀都访问该数组，最后才写回共享内存。遍历
\code{buffer_r} 的循环带有 \code{#pragma unroll}，确保编译器展开循环，使局部
数组始终以常量下标访问，从而可以提升到寄存器。实践中编译器很可能自动展开，显式标注
也能提醒读者这是设计意图。

内核仍先把全局内存数据装入共享内存，再装入寄存器；输出时也先从寄存器写到共享内存，
再写全局内存。共享内存此时不是用于数据复用，而是作为中介，保证全局内存访问合并。
若线程在全局内存与自己的连续寄存器子区段之间直接搬运，warp 内地址会跨步。这与第 6
章的矩阵转角（corner turning）优化相似。

还可继续应用第 6 章的优化清单：装载和存储多个元素时可用向量加载/存储；线程以跨步
模式访问共享内存中的子区段时可能发生 bank 冲突，可通过填充共享内存数组消除。这些优化
留作练习。

\section{内存带宽考量}
\label{sec:scan-memory-bandwidth}

扫描 $N$ 个浮点元素需要从内存加载 $N$ 个浮点数（$4N$ B），执行 $N-1$ 次浮点加法，
再存储 $N$ 个浮点数（$4N$ B），因此算术强度为
\[
  \frac{N-1}{4N+4N}\approx\frac18\ \text{OP/B}.
\]
NVIDIA H100 的峰值浮点吞吐量为 66.9 TFLOPS，峰值全局内存带宽为 3.35 TB/s，计算
至少要达到 $66.9/3.35=20.0$ FLOP/B 才受计算吞吐量限制。扫描远低于这个阈值，是
高度受内存带宽限制的计算。

内存受限计算的峰值性能由峰值带宽决定。H100 上每扫描一个元素至少访问 8 B，所以理论
上限为
\[
  \frac{3.35\times10^{12}}{8}=419\times10^9\ \text{元素/秒}.
\]
达到这个目标要求每个全局内存位置只访问一次，并把计算、同步和共享内存访问的延迟全部
隐藏在全局内存访问之后。实际很难完美利用带宽；即便只在两个数组间复制数据，也不能
达到绝对 100\%。通常，没有冗余全局内存访问且利用率超过 80\% 的内存受限内核，就可
视为充分优化。

图 \ref{fig:11-13} 没有冗余全局内存访问，并且尽力压低了计算、同步与共享内存开销，
因此能够隐藏大部分工作并达到超过 80\% 的带宽利用率。不过，它只让每个线程块局部扫描
自己的区段。把这些区段合并成全局扫描结果还会带来额外全局内存访问和同步开销，这是
下一节要解决的问题。

与归约一样，扫描也由许多库直接提供。C++ 标准库包含
\code{std::inclusive_scan} 与 \code{std::exclusive_scan}，Thrust
\cite{ch11-bell2012thrust,ch11-nvidia-thrust2025} 和 CUB\cite{ch11-nvidia-cub2025} 也提供针对 GPU
高度优化的并行扫描。实际程序通常会使用这些库，但从头实现扫描很适合展示并行化与优化
技术。

\section{合并块区段以完成全局扫描}
\label{sec:global-scan-consolidation}

到目前为止，内核只在输入的各个区段上执行块级扫描。某些应用只需块级结果，但另一些
应用要扫描数百万乃至数十亿个元素，需要多个线程块完成网格级扫描，并合并不同块的区段。

一种方法仍是扫描--扫描--加法，如图 \ref{fig:11-14} 所示：把全局输入划分为各块可以
独立处理的连续区段；每块用图 \ref{fig:11-13} 一类内核扫描自己的区段；各块把末值
（即块内所有元素之和）送入全局块间扫描；最后，每块把此前所有块的总和加到自己的局部
结果上。

\begin{figure}[htbp]
  \centering
  \fbox{\begin{minipage}{0.9\textwidth}\centering
    全局输入划分为连续块区段\\[0.35em]
    $\Downarrow$ 各线程块独立扫描（第一次扫描）\\[0.35em]
    收集块和 $\Downarrow$ 扫描块和（第二次扫描）\\[0.35em]
    各块加上前序块累计和 $\Downarrow$ 全局扫描结果
  \end{minipage}}
  \caption{用扫描--扫描--加法把全局扫描分解为块级扫描}
  \label{fig:11-14}
\end{figure}

主要困难是块间扫描，它要求所有块完成局部扫描后跨块同步。最简单的实现调用三个内核：
第一个扫描块区段，第二个用单个线程块扫描块和，第三个加扫描后的块和。缺点是内核调用
之间必须把中间结果写到全局内存再读回。

设总元素数为 $N$，每块区段大小为 $S$。第一个内核加载并存储各 $N$ 个元素；第二个
内核加载并存储各 $N/S$ 个块和；第三个内核再次加载并存储各 $N$ 个元素。合计访问
\[
 N+N+\frac NS+\frac NS+N+N
 =\left(4+\frac2S\right)N\ \text{个元素},
\]
即 $(16+8/S)N$ B。$S$ 很大时约为每元素 16 B。在 H100 的 3.35 TB/s 带宽下，
峰值只有 $3.35\times10^{12}/16=209\times10^9$ 元素/秒，是理想实现 4190 亿元素/秒
的一半。原因正是第一个内核结束时多写了 $N$ 个中间值，第三个内核开始时又读回它们。

另一种分解是\textbf{归约--扫描--扫描（reduce--scan--scan）}。各块先归约区段，
随后全局扫描块和；最后，每块把此前块的累计和加到区段首元素，再扫描整个块区段，直接
生成最终结果。

\begin{figure}[htbp]
  \centering
  \fbox{\begin{minipage}{0.9\textwidth}\centering
    全局输入划分为连续块区段\\[0.35em]
    $\Downarrow$ 各线程块归约自己的区段\\[0.35em]
    块和 $\Downarrow$ 全局扫描块和\\[0.35em]
    各块把前序累计和放入区段首项并执行块内扫描 $\Downarrow$ 全局结果
  \end{minipage}}
  \caption{用归约--扫描--扫描把全局扫描分解为块级扫描}
  \label{fig:11-15}
\end{figure}

三内核实现中，第一个归约内核加载 $N$ 个元素、存储 $N/S$ 个块和；第二个内核各加载
和存储 $N/S$ 个元素；第三个内核加载 $N$ 个元素并存储 $N$ 个结果。合计访问
\[
 N+\frac NS+\frac NS+\frac NS+N+N
 =\left(3+\frac3S\right)N\ \text{个元素},
\]
即 $(12+12/S)N$ B。$S$ 很大时约为每元素 12 B，对应 H100 上
$279\times10^9$ 元素/秒，是理想值的三分之二，优于扫描--扫描--加法的三内核实现。
它省去了第一个内核对 $N$ 个中间值的存储，但第三个内核仍要重新加载全部输入。

采用归约--扫描--扫描时必须关注运算符的交换律。本章此前的扫描并行化只要求结合律，
而第 10 章某些归约优化还要求交换律。因此，应根据扫描运算符是否可交换，谨慎选择归约
实现。

为何块内按 warp 或线程分解时仍用扫描--扫描--加法？两种分解都有两次扫描，差别是一次
加法与一次归约。归约的优点是写出更少，缺点是归约树需要多个步和中间同步。跨块分解的
首要目标是少访问全局内存，所以归约更有利；块内数据已在寄存器中，多做寄存器加法远比
归约同步便宜，因此扫描--扫描--加法更好。

若要进一步接近理想性能，整个输入必须只加载一次、整个输出只存储一次，多内核方案做
不到这一点。因此，需要在不结束当前内核的情况下完成块间扫描\cite{dotsenko2008fast,
yan2013streamscan,merrill2016single}。

\begin{figure}[htbp]
  \centering\small
  \begin{tabular}{p{0.27\textwidth}|p{0.29\textwidth}|p{0.29\textwidth}}
    (a) 集中式 & (b) 单次回看 & (c) 极端解耦回看\\\hline
    所有块把和交给块 0；块 0 扫描后再分发。需要扫描前后两次网格屏障。
    & 块 $b$ 等待块 $b-1$ 的扫描和，累加自己的和后传给块 $b+1$。
    & 每个块回看并归约此前所有已产生的块和；同步少，但会重复计算。
  \end{tabular}
  \caption{单个内核内实现块间扫描的三种方法}
  \label{fig:11-16}
\end{figure}

图 \ref{fig:11-16}(a) 要求所有块先共享块和，块 0 扫描并重新分发。借助第 18 章将介绍
的协作组可以执行网格级屏障，但它要求所有块同时驻留，因而限制网格中的块数。每块又因
SM 的共享内存与寄存器容量而只能处理有限区段，于是整个可扫描数组也受限。更重要的是，
网格级屏障开销很大，而这里需要两次。

图 \ref{fig:11-16}(b) 的\textbf{单次回看（single lookback）}让块和沿块序列传递并
途中扫描。每个块等待前一块的扫描和，将其与自己的块和相加，再把结果传给下一块。它
不需要网格级屏障，只要求一个块等待此前块完成某项操作。这类同步称为\textbf{单向同步
（unidirectional synchronization）}。它不要求所有块同时活跃，但必须保证较早的块先
被调度。为此，应在线程块开始执行后再动态分配块下标，而不能直接依赖
\code{blockIdx.x}。

图 \ref{fig:11-17} 的设备函数接收块贡献值和动态块下标，并访问扫描和数组及结果
就绪标志数组。块内最后一个
线程拥有块和，作为领导线程。块 0 的前序扫描和显然是
0；其他块反复原子读取前一块的标志，等它就绪后加载其扫描和。

\begin{figure}[htbp]
  \centering
  \begin{minipage}{0.99\textwidth}
  \begin{lstlisting}[language=C++,numbers=left]
__device__ inline float interBlockScan(float val, unsigned int bid,
                                       float *partialSums, unsigned int *flags) {
    __shared__ float prevBlockPartialSum_s;
    if (threadIdx.x == blockDim.x - 1) {
        if (bid > 0) {
            cuda::atomic_ref<unsigned int, cuda::thread_scope_device>
                prevBlockFlag_ref(flags[bid - 1]);
            while (prevBlockFlag_ref.fetch_add(
                       0, cuda::memory_order_acquire) == 0) {}
            prevBlockPartialSum_s = partialSums[bid - 1];
        } else {
            prevBlockPartialSum_s = 0.0f;
        }
        partialSums[bid] = prevBlockPartialSum_s + val;
        cuda::atomic_ref<unsigned int, cuda::thread_scope_device>
            myBlockFlag_ref(flags[bid]);
        myBlockFlag_ref.fetch_add(1, cuda::memory_order_release);
    }
    __syncthreads();
    return prevBlockPartialSum_s;
}
  \end{lstlisting}
  \end{minipage}
  \caption{采用单次回看的块间扫描代码}
  \label{fig:11-17}
\end{figure}

标志必须用设备作用域原子引用访问，才能保证跨块可见性。等待方的原子操作使用
\code{cuda::memory_order_acquire}，表示后续读取块扫描和不能被重排到检查标志之前；
发布方先写 \code{partialSums[bid]}，再以
\code{cuda::memory_order_release} 更新标志，表示此前写内存不能被重排到发布标志之后。
最后的块级屏障保证其他线程在领导线程把前序块扫描和写入共享内存前不会继续。

单次回看虽然能在一个内核中完成扫描，性能却可能很差。它在块之间建立一条很长的依赖链，
成为高延迟关键路径。数组长 $N$、块区段长 $S$ 时有 $N/S$ 个块；集中式方案需要两次
网格屏障，单次回看却需要 $N/S$ 次单向同步。单向同步虽轻量，数量过多仍会使算法低效。

\textbf{多次回看（multiple lookback）}允许一个块查看多个前序块，从而缩短关键路径。
例如块 $b$ 发现 $b-1$ 只有块和、尚无扫描和时，可继续查看 $b-2$。若 $b-2$ 的扫描和
已就绪，$b$ 就能自己把 $b-1$ 的块和加上去，而不必等 $b-1$ 完成。这样，$b$ 和
$b-1$ 都会做同一次加法，增加了冗余工作，却减少了等待延迟。

回看范围可以固定，也可以持续向前直到找到一个已就绪的扫描和。后一种称为
\textbf{解耦回看（decoupled lookback）}。回看越远，重复工作越多，但绕过的依赖链也
越长。图 \ref{fig:11-16}(c) 是极端情形：每个块归约此前所有块和；它以更多计算换来
更少的步和同步，类似第 \ref{sec:kogge-stone-scan} 节的朴素并行扫描。块间同步远比计算
昂贵，而且块和数远少于总元素数，所以此处较差的理论工作效率反而可能更有利。

图 \ref{fig:11-18} 给出单内核全局扫描。开头动态分配块下标：块开始执行后，线程 0
原子递增块计数器取得下标，再通过共享内存广播给块内线程。如此，实际先调度的块得到较小
下标，满足单向同步要求。随后代码与图 \ref{fig:11-13} 相同：装载块区段，各线程扫描
自己的寄存器子区段，扫描线程和，并加上前序线程累计和。接着调用图
\ref{fig:11-17} 的块间扫描，把前序块累计和加到所有局部结果，最后合并写回。

\begin{figure}[htbp]
  \centering
  \begin{minipage}{0.99\textwidth}
  \begin{lstlisting}[language=C++,numbers=left]
__global__ void scan_kernel(float *input, float *output,
                            unsigned int *blockCounter, float *partialSums,
                            unsigned int *flags, unsigned int N) {
    __shared__ unsigned int bid_s;
    if (threadIdx.x == 0) {
        cuda::atomic_ref<unsigned int, cuda::thread_scope_device>
            blockCounter_ref(*blockCounter);
        bid_s = blockCounter_ref.fetch_add(1, cuda::memory_order_relaxed);
    }
    __syncthreads();
    unsigned int bid = bid_s;
    unsigned int blockSegment = COARSE_FACTOR * bid * blockDim.x;

    __shared__ float buffer_s[COARSE_FACTOR * BLOCK_DIM];
    for (unsigned int c = 0; c < COARSE_FACTOR; ++c)
        buffer_s[c * BLOCK_DIM + threadIdx.x] =
            input[blockSegment + c * BLOCK_DIM + threadIdx.x];
    __syncthreads();

    unsigned int threadSegment = threadIdx.x * COARSE_FACTOR;
    float buffer_r[COARSE_FACTOR];
    buffer_r[0] = buffer_s[threadSegment];
#pragma unroll
    for (unsigned int c = 1; c < COARSE_FACTOR; ++c)
        buffer_r[c] = buffer_s[threadSegment + c] + buffer_r[c - 1];

    float threadSum = blockScan(buffer_r[COARSE_FACTOR - 1]);
    __shared__ float threadSums[BLOCK_DIM];
    threadSums[threadIdx.x] = threadSum;
    __syncthreads();
    float prevPartialSum = (threadIdx.x > 0)
                           ? threadSums[threadIdx.x - 1] : 0.0f;
#pragma unroll
    for (unsigned int c = 0; c < COARSE_FACTOR; ++c)
        buffer_r[c] += prevPartialSum;

    float prevBlockPartialSum = interBlockScan(
        buffer_r[COARSE_FACTOR - 1], bid, partialSums, flags);
#pragma unroll
    for (unsigned int c = 0; c < COARSE_FACTOR; ++c)
        buffer_s[threadSegment + c] = buffer_r[c] + prevBlockPartialSum;
    __syncthreads();

    for (unsigned int c = 0; c < COARSE_FACTOR; ++c)
        output[blockSegment + c * BLOCK_DIM + threadIdx.x] =
            buffer_s[c * BLOCK_DIM + threadIdx.x];
}
  \end{lstlisting}
  \end{minipage}
  \caption{通过块间扫描合并区段的单内核扫描}
  \label{fig:11-18}
\end{figure}

\section{使用 Brent--Kung 算法进行并行扫描}
\label{sec:brent-kung-scan}

本章此前一直以 Kogge--Stone 为基础，它速度快、概念简单，却工作效率较差。本节介绍
Brent--Kung 算法：它以更多计算步骤为代价，减少总工作量。线程粗化已经能缓解
Kogge--Stone 的工作效率问题，但理解另一种取舍仍有价值。

观察 Kogge--Stone 的执行过程，会发现一些中间结果还可进一步共享。不过，为了让更多
线程共享，必须有策略地计算并分发中间值，因而需要额外步骤。用二元运算符对 $N$ 个值
求和的最快并行方法是归约树，在执行单元充足时用 $\log_2N$ 个时间单位得到总和，同时
产生一些可供扫描输出复用的部分和。这个观察既是 Kogge--Stone 加法器的基础，也是
Brent--Kung 加法器\cite{brent1982regular} 的基础。

图 \ref{fig:11-19} 展示 16 元素的 Brent--Kung 包含式扫描。上半部是归约树：第一步
只更新下标形如 $2n-1$ 的位置，第二步更新 $4n-1$，第三步更新 $8n-1$，第四步只更新
$16n-1=15$。它用 $8+4+2+1=15$ 次操作得到总和。一般地，归约阶段完成
\[
 \frac N2+\frac N4+\cdots+2+1=N-1
\]
次操作。

\begin{figure}[htbp]
  \centering\scriptsize
  \begin{tabular}{c|l}
    归约距离 1 & 更新 $1,3,5,7,9,11,13,15$\\
    归约距离 2 & 更新 $3,7,11,15$\\
    归约距离 4 & 更新 $7,15$\\
    归约距离 8 & 更新 $15$，得到 $\sum_{0}^{15}x_i$\\\hline
    反向距离 4 & 用位置 7 补全位置 11\\
    反向距离 2 & 用位置 3、7、11 补全位置 5、9、13\\
    反向距离 1 & 补全位置 $2,4,6,8,10,12,14$
  \end{tabular}
  \caption{基于 Brent--Kung 加法器设计的并行包含式扫描}
  \label{fig:11-19}
\end{figure}

算法下半部用反向树把部分和分发给需要它们的位置。归约树中的每次加法都累积一个连续
范围，所以任一位置的当前值都可写成 $x_i\ldots x_j$。归约结束后，位置 0、1、3、7、
15 已经是最终扫描值，其余位置还需补入左侧的若干连续部分和。

\begin{figure}[htbp]
  \centering\scriptsize
  \begin{tabular}{c|cccccccc}
    位置 & 0 & 1 & 2 & 3 & 5 & 7 & 11 & 15\\\hline
    归约后 & $x_0$ & $x_0..x_1$ & $x_2$ & $x_0..x_3$ & $x_4..x_5$
      & $x_0..x_7$ & $x_8..x_{11}$ & $x_0..x_{15}$\\
    反向树后 & $x_0$ & $x_0..x_1$ & $x_0..x_2$ & $x_0..x_3$ & $x_0..x_5$
      & $x_0..x_7$ & $x_0..x_{11}$ & $x_0..x_{15}$
  \end{tabular}
  \caption{反向树各层之后数组中部分和的演化（概念示意）}
  \label{fig:11-20}
\end{figure}

图 \ref{fig:11-20} 中，一个位置最多需要来自其他位置的 $\log_2N-1$ 个部分和，这些
来源之间的距离依次是 $1,2,4,\ldots$。例如 $N=16$ 时，位置 14 最初只有 $x_{14}$，
还需要位置 13 的 $x_{12}..x_{13}$、位置 11 的 $x_8..x_{11}$ 和位置 7 的
$x_0..x_7$。反向树先处理相距 4 的来源，再处理相距 2 和 1 的来源：第一层补全位置
11；第二层补全 5、9、13；最后一层补全 2、4、6、8、10、12、14。

反向树的运算数为
\[
 \left(\frac N8-1\right)+\left(\frac N4-1\right)
 +\left(\frac N2-1\right)=N-1-\log_2N.
\]
加上归约阶段的 $N-1$ 次，总运算数是
$2N-2-\log_2N=O(N)$，而 Kogge--Stone 为 $O(N\log N)$。

两种算法的取舍很清楚：Brent--Kung 工作效率更高、总操作更少，但执行的步数更多。不过，
应用线程粗化后，扫描的大部分工作已经由工作高效的线程级顺序扫描完成，工作低效的并行
扫描只发生在 warp 级，只占很小一部分运行时间。

实践中，在 warp 级 Kogge--Stone 通常反而更快。SIMD 执行使 Brent--Kung 避免的工作
变成不活跃 lane，这些 lane 仍消耗执行资源；与其让它们空闲，不如像 Kogge--Stone 那样
多做一些工作而用更少步骤结束。尽管如此，Brent--Kung 仍是理解理论工作效率与实际实现
因素相互作用的有趣案例，其 CUDA 实现留作练习。

\section{小结}
\label{sec:scan-summary}

本章研究了并行扫描（前缀和）这一重要并行模式。它能为需求量不同的参与方并行分配资源，
把基于数学递推、看似顺序的计算转换为并行计算，从而消除许多应用中的顺序瓶颈。简单顺序
扫描对 $N$ 个元素只做 $N-1=O(N)$ 次加法。

首先介绍的 Kogge--Stone 算法快速而直观，但并非工作高效：它完成
$O(N\log_2N)$ 次操作，多于顺序算法。数据规模越大，并行算法要胜过顺序算法所需的执行
单元也越多。因此，它通常用在执行资源丰富的处理器上，扫描大小适中的区段。

我们用双缓冲和 warp 级原语减少同步开销，又用线程粗化提高工作效率：每个线程先对自己的
输入子区段执行工作高效的顺序扫描，再由块内线程协作完成工作效率较低的块级并行扫描。

为了跨多个线程块合并区段，本章比较了扫描--扫描--加法与归约--扫描--扫描。三内核实现中，
后者的冗余全局内存访问更少；单内核扫描还可以通过块间单向同步进一步减少全局内存访问。
单向同步必须配合动态块下标分配以避免死锁；解耦回看则用于缩短它建立的关键依赖路径。

最后，Brent--Kung 用归约树和反向树把工作量降到 $O(N)$，但需要更多步骤。线程粗化与
warp 的 SIMD 特性会削弱这一理论优势。GPU 扫描实现与优化相当复杂，实际编程更推荐使用
Thrust 或 CUB，而不是从头编写；但扫描仍是研究并行模式优化取舍的重要案例。

\phantomsection
\section*{练习}
\addcontentsline{toc}{section}{练习}

\begin{enumerate}[label=\arabic*.,leftmargin=2.7em]
  \item 对数组 $[4,6,7,1,2,8,5,2]$ 使用 Kogge--Stone 算法执行并行包含式前缀扫描，
  写出每一步后的数组中间状态。

  \item 分析图 \ref{fig:11-3} 的 Kogge--Stone 内核，证明步幅不超过 warp 大小一半时，
  控制分歧只发生在每个块的第一个 warp 中。也就是说，warp 大小为 32 时，步幅为
  1、2、4、8、16 的各轮会出现控制分歧。

  \item 图 \ref{fig:11-3} 的 Kogge--Stone 内核扫描 2048 个元素时，总共执行多少次
  加法？

  \item 修改图 \ref{fig:11-13} 的内核，使用向量加载与存储，并消除共享内存 bank 冲突。

  \item 修改图 \ref{fig:11-17} 的 \code{interBlockScan} 设备函数，用解耦回看替代
  单次回看。

  \item 对数组 $[4,6,7,1,2,8,5,2]$ 使用 Brent--Kung 算法执行并行包含式前缀扫描，
  写出每一步后的数组中间状态。

  \item 实现一个在块级扫描数组的 Brent--Kung CUDA 内核。
\end{enumerate}

\begin{chapterbibliography}{9}
  \bibitem{ch11-bell2012thrust}
  N. Bell, J. Hoberock, Thrust: a productivity-oriented library for CUDA,
  in: \emph{GPU Computing Gems Jade Edition}, Elsevier, 2012, pp. 359--371.

  \bibitem{ch11-nvidia-thrust2025}
  NVIDIA, Thrust: the C++ parallel algorithms library, 2025.
  \url{https://nvidia.github.io/cccl/thrust/index.html}.

  \bibitem{ch11-nvidia-cub2025}
  NVIDIA, CUB, 2025.
  \url{https://nvidia.github.io/cccl/cub/index.html}.

  \bibitem{kogge1973parallel}
  P. M. Kogge, H. S. Stone, A parallel algorithm for the efficient solution of a
  general class of recurrence equations, \emph{IEEE Transactions on Computers}
  100(8) (1973) 786--793.

  \bibitem{dotsenko2008fast}
  Y. Dotsenko, N. K. Govindaraju, P.-P. Sloan, C. Boyd, J. Manferdelli,
  Fast scan algorithms on graphics processors, in: \emph{Proceedings of the 22nd
  Annual International Conference on Supercomputing}, 2008, pp. 205--213.

  \bibitem{yan2013streamscan}
  S. Yan, G. Long, Y. Zhang, StreamScan: fast scan algorithms for GPUs without
  global barrier synchronization, in: \emph{Proceedings of the 18th ACM SIGPLAN
  Symposium on Principles and Practice of Parallel Programming}, 2013, pp. 229--238.

  \bibitem{merrill2016single}
  D. Merrill, M. Garland, Single-pass parallel prefix scan with decoupled look-back,
  Technical Report NVR-2016-002, NVIDIA, 2016.

  \bibitem{brent1982regular}
  R. P. Brent, H. T. Kung, A regular layout for parallel adders,
  \emph{IEEE Transactions on Computers} 100(3) (1982) 260--264.
\end{chapterbibliography}
